\chapter{Введение: почему программа меняет основной объект}
\label{ch:v8-introduction-and-problem-statement}

\section{Что требовалось после Тома VII}

Том VII получил устойчивый качественный переход
\begin{equation}
 (7,0,20)_0\longrightarrow(0,0,27)_*,
\end{equation}
но не вывел единственную массовую метрику, общий физический след и семейные
фазы. Его заключение оставило три возможных входа: межсекторный коннектор,
второй семейный тензор или общий калибровочно-пространственный источник нормировки.

Перед открытием нового тома все три входа были проверены повторно. Прямой
Коннектор Мориты потребовал невыведенного $U(11)$-отождествления; единственный
совместимый с вещественной структурой цветной канал был уничтожен бесцветным проектором; активный
Ходж-вакуум имел нулевой калибровочный индекс; старые семейные тензоры либо
предполагали уже закрытый коннектор, либо провалили слепой тест, либо
оставались невыбранным меню. Производная проектора ядра дала каноническую
кинематику фиксированного ранга, но расходилась на общем пути изменения ранга.

Следовательно, Том VIII нельзя строить как ещё одно продолжение прежнего
поля с новым коэффициентом.

\section{Пропущенная ранняя предпосылка}

Археология исходных TOE--UGSM-текстов показала, что ранним фундаментальным
объектом был полный нелокальный корреляционный оператор $C(x,y)$, тогда как
поздняя программа в основном использовала его спектр, тепловой след и
конечные моменты. Эти редукции уничтожают собственные направления и
двухточечную локальность.

На старой коспектральной паре Тома V полный оператор
\begin{equation}
 C_\tau=e^{-\tau H}
\end{equation}
точно различил две геометрии, неразличимые по всем скалярным спектральным
функционалам, а формула
\begin{equation}
 H=-\tau^{-1}\log C_\tau
\end{equation}
вернула их конечную связность. Тем самым был найден промежуточный объект:
полное ядро вместе с алгеброй наблюдаемых, но без предположения готовой
спектральной геометрии.

\section{Почему том разрешено открыть}

Проект уже содержит не только абстрактную возможность такого объекта.
Физическая инцидентность $A_0:\mathbb C^{11}\to\mathbb C^{10}$ задаёт
связывающий оператор
\begin{equation}
 D_A=\begin{pmatrix}0&A_0^*\\A_0&0\end{pmatrix}
\end{equation}
и коэффициентно-свободный квантово-марковский генератор
\begin{equation}
 \mathcal L_A=-\frac12\operatorname{ad}_{D_A}^2
\end{equation}
на концевой алгебре $M_{11}(\mathbb C)\oplus M_{10}(\mathbb C)$.
Полугруппа вполне положительна, сохраняет единицу и след и имеет ненулевой
перенос между двумя концевыми углами.

Полный калибровочный набор сокращает её 41-мерную неподвижную алгебру до
\begin{equation}
 \mathcal Z_{\rm fix}
 =\mathbb CP_q\oplus\mathbb CP_\ell\simeq\mathbb C^2,
 \qquad (\rank P_q,\rank P_\ell)=(12,9).
\end{equation}
Это первый новый примитив: не выбранный коннектор и не подогнанная масса, а
выведенный операторный процесс с вычислимой центральной границей.

\section{Центральная задача Тома VIII}

Том проверяет, может ли полное корреляционное ядро и его составная
кривизна породить калибровочно-совместимый переход между $P_q$ и $P_\ell$, не
создавая цветного вакуумного среднего и не вводя произвольное
отождествление пространств.

Рабочая гипотеза состоит в том, что межсекторная связь может возникнуть не
как линейная физическая стрелка, а как цветосохраняющий составной канал
полного ядра. Нулевая гипотеза состоит в том, что $\mathbb C^2$ является
строгой суперселекционной границей текущего физического носителя.

\section{Порядок и стоп-критерий}

\begin{enumerate}
 \item Классифицировать все составные блоки полного ядра между $P_q$ и
 $P_\ell$ с учётом калибровочного условия, вещественной структуры и условия Ходжа.
 \item Проверить, возникает ли ненулевой межсекторный блок без введения
 нового линейного коннектора.
 \item При положительном результате построить одну полугруппу и проверить
 её неподвижную алгебру, полный гессиан и ограниченность действия.
 \item При отрицательном результате доказать суперселекцию и закрыть
 текущий конечный носитель до расширения алгебры или поля.
\end{enumerate}

\begin{nogocriterion}[Стоп-критерий Тома VIII]
Если вся калибровочно-инвариантная операторная оболочка полного ядра сохраняет
$P_q$ и $P_\ell$ как центральные проекторы, том не продолжает поиск масс и
смешиваний. Тогда двухсекторность объявляется строгой границей данной
архитектуры, а не почти полученным объединением.
\end{nogocriterion}

\begin{protocol}[Открытие Тома VIII]
\begin{enumerate}
 \item ранние предпосылки перечитаны: \textbf{да};
 \item старые три входа повторно проверены: \textbf{закрыты};
 \item новый фундаментальный объект: \textbf{полное корреляционное ядро};
 \item алгебра наблюдаемых: \textbf{$M_{11}\oplus M_{10}$};
 \item квантово-марковская полугруппа: \textbf{выведена};
 \item классический центральный остаток: \textbf{$\mathbb C^2$};
 \item межсекторный переход: \textbf{не утверждается};
 \item центральный вычислимый тест: \textbf{сохранение или снятие
 $P_q,P_\ell$};
 \item феноменологическое чтение: \textbf{запрещено до прохождения теста}.
\end{enumerate}
\end{protocol}